\chapter{Grupy Liego i ich algebry}
\label{ch:grupy-liego}
\index{grupa Liego}
\index{algebra Liego}

W Rozdziale~\ref{ch:rozmaitosci} grupa ortogonalna pojawiła się jako
rozmaitość, a w Rozdziale~\ref{ch:pola} poznaliśmy przepływy i nawias pól.
Teraz połączymy te konstrukcje. W grupie Liego każdy wektor styczny
przy elemencie neutralnym wyznacza pole na całej grupie, a przepływ
tego pola daje rodzinę symetrii zależną od jednego parametru.
Grupy i homomorfizmy wystąpiły już w
Definicji~\ref{def:grupa-hom-5}; tutaj wymagamy ponadto gładkości.

\section{Grupa z gładką strukturą}
\label{sec:grupy-liego-przyklady}

\begin{definicja}[Grupa Liego]
\label{def:grupa-liego}
\index{grupa Liego!definicja}
\emph{Grupą Liego} nazywamy skończenie wymiarową gładką rozmaitość
bez brzegu $G$, wyposażoną w strukturę grupy, dla której
\[
 m:G\times G\longrightarrow G,\quad m(g,h)=gh,
 \qquad \iota:G\longrightarrow G,\quad\iota(g)=g^{-1}
\]
są gładkie. Jej element neutralny oznaczamy $e$.
Intuicyjnie można w niej wykonywać działania grupowe i jednocześnie
różniczkować je względem parametrów. Homomorfizm grup Liego nazywamy
\emph{gładkim}, gdy jest gładki jako odwzorowanie rozmaitości.
\end{definicja}

Dla $g\in G$ oznaczamy przez $L_g(h)=gh$ i $R_g(h)=hg$
translacje. Obie są dyfeomorfizmami: są gładkie jako ograniczenia
$m$, a ich odwrotnościami są odpowiednio $L_{g^{-1}}$ i
$R_{g^{-1}}$. Pozwalają przenosić rachunek z punktu $e$ do
każdego punktu grupy.

\begin{przyklad}[Grupy macierzy odwracalnych]
\label{prz:gl-liego}
\index{grupa ogólna liniowa}
Grupa $\mathrm{GL}(n,\R)$ z
Dodatku~\ref{sec:app-grupa-gl} jest otwartym podzbiorem
$\R^{n^2}$, bo wyznacznik jest wielomianem.
Iloczyn macierzy jest wielomianowy w ich współczynnikach, a
\[
 A^{-1}=\frac{\operatorname{adj}A}{\det A}
\]
jest gładki na tym otwartym zbiorze. Stąd $\mathrm{GL}(n,\R)$
jest grupą Liego wymiaru $n^2$. Ten sam argument, po potraktowaniu
części rzeczywistych i urojonych współczynników jako współrzędnych,
daje $\mathrm{GL}(n,\C)$ wymiaru rzeczywistego $2n^2$.
W obu przypadkach $T_I\mathrm{GL}(n,\mathbb K)=M_n(\mathbb K)$
jako przestrzeń rzeczywista: otwarty podzbiór przestrzeni macierzy
ma tę samą przestrzeń styczną co otoczenie.
\end{przyklad}

\begin{przyklad}[Grupy ortogonalne]
\label{prz:orth-liego}
\index{grupa ortogonalna}
W Przykładzie o grupie ortogonalnej w
Rozdziale~\ref{ch:rozmaitosci} wykazaliśmy, że
grupa $\mathrm O(n)$ z Dodatku~\ref{sec:app-grupa-o}
jest gładką rozmaitością i że
$T_I\mathrm O(n)=\{X:X^T=-X\}$. Mnożenie jest ograniczeniem
wielomianowego mnożenia macierzy, a odwrotność ma tutaj prostą
postać $A^{-1}=A^T$, więc oba działania są gładkie.
Ponieważ $\det A\in\{-1,1\}$ na $\mathrm O(n)$,
$\mathrm{SO}(n)$ jest
otwartą i domkniętą podrozmaitością $\mathrm O(n)$.
Jest podgrupą Liego tego samego wymiaru $n(n-1)/2$.
Jej przestrzeń styczna przy $I$ również składa się z
macierzy skośnie symetrycznych.
\end{przyklad}

\begin{przyklad}[Torus jako grupa]
\label{prz:torus-liego}
\index{torus!jako grupa Liego}
Na torusie $T^k=\R^k/\Z^k$ z Przykładu~\ref{prz:torus}
kładziemy $[x]+[y]=[x+y]$ i $-[x]=[-x]$. W lokalnych mapach
utworzonych z rzutowania $\R^k\to T^k$ działania mają postać
zwykłego dodawania i zmiany znaku, z ewentualną całkowitą
translacją współrzędnych. Są więc gładkie. Równoważny model
\[
 T^k\longrightarrow(S^1)^k,\qquad
 [x_1,\ldots,x_k]\longmapsto
 (e^{2\pi i x_1},\ldots,e^{2\pi i x_k})
\]
przekształca dodawanie w mnożenie zespolone na każdej współrzędnej;
odwrotność staje się sprzężeniem zespolonym. Dla $k=2$ torus
obrotowy z Rozdziału~\ref{ch:rozmaitosci} otrzymuje tę strukturę
po przeniesieniu jej przez gładką parametryzację
$(u,v)\mapsto((R+r\cos v)\cos u,(R+r\cos v)\sin u,r\sin v)$,
gdzie $R>r>0$. Jego mnożenie zależy od tej identyfikacji, a nie
od dodawania punktów w otaczającej go przestrzeni $\R^3$.
\end{przyklad}

\begin{przyklad}[Macierze specjalnie unitarne]
\label{prz:su2-liego}
\index{grupa specjalna unitarna}
Niech $\mathrm{SU}(2)$ będzie grupą z
Dodatku~\ref{sec:app-grupa-u}.
Z warunków $U^{-1}=U^*$ i $\det U=1$ wynika, przez porównanie
wzorów na odwrotność, że każdy jej element ma jedyną postać
\begin{equation}\label{eq:su2-macierz}
 U(a,b)=\begin{pmatrix}a&b\\-\bar b&\bar a\end{pmatrix},
 \qquad |a|^2+|b|^2=1.
\end{equation}
Odwrotnie, bezpośrednie pomnożenie pokazuje
$U(a,b)^*U(a,b)=I$ i $\det U(a,b)=|a|^2+|b|^2=1$.
Współrzędne $(\Re a,\Im a,\Re b,\Im b)$ identyfikują zatem
$\mathrm{SU}(2)$ gładko ze sferą $S^3\subset\R^4$.
Mnożenie macierzy ogranicza się do gładkiego odwzorowania
$S^3\times S^3\to S^3$, a odwrotność $U^{-1}=U^*$ jest
ograniczeniem rzeczywiście liniowego odwzorowania na $M_2(\C)$.
Jest to więc zwarta, trójwymiarowa grupa Liego.
\end{przyklad}

\section{Od wektora przy jedności do algebry Liego}
\label{sec:algebra-liego}

\begin{definicja}[Pole lewostronnie niezmiennicze i algebra Liego]
\label{def:algebra-liego}
\index{pole lewostronnie niezmiennicze}
\index{algebra Liego!definicja}
Pole $V$ na $G$ jest \emph{lewostronnie niezmiennicze}, gdy
$(L_g)_*V=V$ dla każdego $g\in G$. Dla $X\in T_eG$ oznaczamy
\[
 X^L_g:=(\dd L_g)_eX.
\]
\emph{Algebrą Liego} grupy $G$ nazywamy przestrzeń
$\mathfrak g=T_eG$ z nawiasem
\[
 [X,Y]_{\mathfrak g}:=[X^L,Y^L]_e.
\]
Wektor przy jedności opisuje ruch nieskończenie mały; translacje
przenoszą ten sam ruch do wszystkich punktów. Przykładami są
$\mathfrak{gl}(n,\R)=M_n(\R)$ i, jak pokażemy niżej,
$\mathfrak t^k\cong\R^k$ z zerowym nawiasem.
\end{definicja}

\begin{twierdzenie}[Niezmienniczość i zamknięcie na nawias]
\label{thm:lie-algebra-fields}
Odwzorowanie $X\mapsto X^L$ jest izomorfizmem liniowym
$T_eG$ na przestrzeń pól lewostronnie niezmienniczych.
Nawias dwóch takich pól jest znowu lewostronnie niezmienniczy.
W szczególności $\mathfrak g$ jest algebrą Liego: jej nawias jest
dwuliniowy, antysymetryczny i spełnia tożsamość Jacobiego.
\end{twierdzenie}
\begin{proof}
Gładkość $X^L$ wynika z gładkości różniczki mnożenia po
ograniczeniu do $G\times\{e\}$. Dla $g,h\in G$ mamy
$L_g\circ L_h=L_{gh}$, więc reguła łańcucha daje
\[
 (\dd L_g)_h X^L_h
 = (\dd L_g)_h(\dd L_h)_eX
 = (\dd L_{gh})_eX=X^L_{gh}.
\]
Jest to niezmienniczość. Jeśli $V$ jest dowolnym polem
lewostronnie niezmienniczym, to po podstawieniu $h=e$
otrzymujemy $V_g=(\dd L_g)_eV_e$; zatem $V=(V_e)^L$.
Wartość przy $e$ jest więc odwrotnością odwzorowania
$X\mapsto X^L$, a obie mapy są liniowe.

Dla dyfeomorfizmu $F$ zachodzi
$F_*[V,W]=[F_*V,F_*W]$. Można to sprawdzić bezpośrednio
na funkcji $f$: definicja pola przeniesionego daje
$(F_*V)(f)=(V(f\circ F))\circ F^{-1}$, a po dwukrotnym
zastosowaniu tego wzoru różnica złożeń staje się
$([V,W](f\circ F))\circ F^{-1}$.
Stosując tę równość do $F=L_g$ i $V=X^L$, $W=Y^L$,
dostajemy
\[
 (L_g)_*[X^L,Y^L]
 =[(L_g)_*X^L,(L_g)_*Y^L]=[X^L,Y^L].
\]
Własności nawiasu pól zostały dowiedzione w
Fakcie~\ref{fakt:nawias-wlasnosci}. Przez izomorfizm
$X\mapsto X^L$ przechodzą na nawias w $T_eG$.
\end{proof}

\begin{przyklad}[Komutator macierzy]
\label{prz:lie-bracket-matrix}
\index{komutator macierzy}
Na otwartym $\mathrm{GL}(n,\R)$ pole odpowiadające macierzy
$X$ ma wartość $X^L_A=AX$. Dla $Y^L_A=AY$ różniczka
odwzorowania $A\mapsto AY$ w kierunku $H$ jest równa $HY$.
Wzór współrzędny na nawias z~\eqref{eq:nawias-wspolrzedne}
daje więc
\[
 [X^L,Y^L]_A
 =\D(Y^L)_A[AX]-\D(X^L)_A[AY]
 =AXY-AYX=A(XY-YX).
\]
Stąd $[X,Y]_{\mathfrak{gl}}=XY-YX$. To samo obliczenie,
po ograniczeniu pól stycznych do podgrupy macierzy, obowiązuje
dla $\mathrm O(n)$, $\mathrm{SO}(n)$ i $\mathrm{SU}(2)$.
W szczególności
$\mathfrak{so}(n)=\{X:X^T=-X\}$ jest zamknięta na komutator,
bo $(XY-YX)^T=-(XY-YX)$.

Dla $\mathrm{SU}(2)$ różniczkowanie $U^*U=I$ i $\det U=1$
w punkcie $I$ daje odpowiednio $X^*+X=0$ i
$\operatorname{tr}X=0$; różniczka wyznacznika w kierunku
$X$ wynosi $\operatorname{tr}X$, co widać z rozwinięcia
$\det(I+tX)=1+t\operatorname{tr}X+t^2\det X$.
Ponieważ obie strony mają wymiar
rzeczywisty trzy, otrzymujemy
\[
 \mathfrak{su}(2)
 =\left\{\begin{pmatrix}i\alpha&z\\-\bar z&-i\alpha\end{pmatrix}
 :\alpha\in\R,\ z\in\C\right\}.
\]
Na torusie translacje w mapach ilorazowych mają stałą
różniczkę równą identyczności. Odpowiadające im pola mają
stałe współczynniki, więc ich nawias znika:
$\mathfrak t^k\cong(\R^k,0)$.
\end{przyklad}

\section{Podgrupy jednoparametrowe i eksponenta}
\label{sec:lie-exponential}

\begin{definicja}[Podgrupa jednoparametrowa]
\label{def:one-parameter}
\index{podgrupa jednoparametrowa}
\emph{Podgrupą jednoparametrową} nazywamy gładki
homomorfizm $\gamma:(\R,+)\to G$.
Warunek homomorfizmu ma postać
$\gamma(s+t)=\gamma(s)\gamma(t)$ i $\gamma(0)=e$.
Na $S^1$ przykładem jest $t\mapsto e^{it}$; ruch wraca
do jedności po czasie $2\pi$, więc homomorfizm nie musi być
różnowartościowy.
\end{definicja}

\begin{twierdzenie}[Wektor styczny wyznacza jedną podgrupę]
\label{thm:one-param-lie}
Dla każdego $X\in\mathfrak g$ istnieje dokładnie jedna
podgrupa jednoparametrowa $\gamma_X$ taka, że
$\gamma_X'(0)=X$. Jest ona krzywą całkową pola $X^L$
przez $e$ i zachodzi
\[
 \Phi_t^{X^L}(g)=g\gamma_X(t)
 \qquad (g\in G,\ t\in\R).
\]
\end{twierdzenie}
\begin{proof}
Twierdzenie~\ref{thm:przeplyw} daje lokalną krzywą
całkową $\gamma$ pola $X^L$ z $\gamma(0)=e$.
Jeżeli $\eta(t)$ jest krzywą całkową przez $h$, to
$L_{h^{-1}}\eta(t)$ również jest krzywą całkową $X^L$:
niezmienniczość pola i reguła łańcucha przenoszą równanie
$\eta'=X^L_\eta$ na tę krzywą. Jednoznaczność rozwiązań
z Rozdziału~\ref{ch:pola} daje zatem, dla małych $t$,
$\Phi_t^{X^L}(h)=h\gamma(t)$.

Wykażmy, że pole jest zupełne. Wybierzmy $\varepsilon>0$,
tak aby $\gamma(t)$ istniała dla $|t|<\varepsilon$.
Z niezmienniczości $X^L$ krzywa
$s\mapsto h\gamma(s)$, $|s|<\varepsilon$, jest
rozwiązaniem przez \emph{dowolny} $h\in G$ przez taki sam
przedział czasu. Gdyby maksymalna krzywa przez $e$ miała
skończony prawy koniec $b$, wybierając $t_0<b$ z
$b-t_0<\varepsilon/2$, przedłużylibyśmy ją wzorem
$\gamma(t_0+s)=\gamma(t_0)\gamma(s)$ do
$t_0+s>b$, wbrew maksymalności. Lewy koniec traktujemy
tak samo. Stąd $\gamma$ istnieje dla wszystkich $t$.

Dla ustalonego $s$ krzywa $t\mapsto\gamma(s)\gamma(t)$
jest całkowa i w chwili $t=0$ przechodzi przez
$\gamma(s)$. Krzywa $t\mapsto\gamma(s+t)$ ma te same
własności. Jednoznaczność daje
$\gamma(s+t)=\gamma(s)\gamma(t)$, zatem $\gamma$ jest
podgrupą jednoparametrową. Ponownie przez translację
jej przepływ przez $g$ ma podaną postać.

Jeśli $\beta:\R\to G$ jest inną gładką podgrupą
jednoparametrową z $\beta'(0)=X$, to
\[
 \beta'(t)=\left.\frac{\dd}{\dd s}\right|_{s=0}
 \beta(t)\beta(s)=(\dd L_{\beta(t)})_eX=X^L_{\beta(t)}.
\]
Jest więc krzywą całkową tego samego pola przez $e$;
jednoznaczność przepływu wymusza $\beta=\gamma$.
\end{proof}

\begin{definicja}[Odwzorowanie wykładnicze]
\label{def:lie-exp}
\index{odwzorowanie wykładnicze!grupy Liego}
\emph{Odwzorowaniem wykładniczym} grupy Liego nazywamy
\[
 \exp_G:\mathfrak g\longrightarrow G,\qquad
 \exp_G(X)=\gamma_X(1).
\]
Wektor $X$ określa prędkość przy jedności, a $\exp_G(X)$
jest punktem osiągniętym po czasie jeden przy tej stałej,
przenoszonej translacjami prędkości. Dla torusa z
Przykładu~\ref{prz:torus-liego}, przy identyfikacji
$T_eT^k\cong\R^k$ z mapy ilorazowej, mamy
$\exp_{T^k}(v)=[v]$.
\end{definicja}

\begin{stwierdzenie}[Własności eksponenty]
\label{prop:lie-exp-properties}
Odwzorowanie $\exp_G$ jest gładkie,
$(\dd\exp_G)_0=\id_{\mathfrak g}$ oraz
\[
 \exp_G((s+t)X)=\exp_G(sX)\exp_G(tX),\qquad
 \exp_G(-X)=\exp_G(X)^{-1}.
\]
W szczególności $\exp_G$ jest dyfeomorfizmem pewnego
otoczenia zera na otoczenie $e$.
\end{stwierdzenie}
\begin{proof}
Liniowa zależność $X\mapsto X^L$ i gładkość różniczki
mnożenia sprawiają, że $(X,g)\mapsto X^L_g$ jest
gładkie. Na rozmaitości $\mathfrak g\times G$ rozważmy
pole $(X,g)\mapsto(0,X^L_g)$. Jego przepływ ma postać
$(X,g)\mapsto(X,\Phi_t^{X^L}(g))$ i jest określony
dla wszystkich czasów na mocy
Twierdzenia~\ref{thm:one-param-lie}. Gładkość przepływu
z Twierdzenia~\ref{thm:przeplyw} daje więc gładkość
$(X,t)\mapsto\gamma_X(t)$, a w szczególności $\exp_G$.
Jednoznaczność podgrupy o zadanej pochodnej
daje $\gamma_X(t)=\exp_G(tX)$: dla każdej liczby $a$
krzywa $s\mapsto\gamma_X(as)$ ma pochodną $aX$ w zerze,
więc $\gamma_{aX}(1)=\gamma_X(a)$.
Prawo grupowe z poprzedniego dowodu daje pierwszą równość,
a podstawienie $s=-t$ daje drugą.

Dla $X\in\mathfrak g$ mamy
\[
 (\dd\exp_G)_0(X)
 =\left.\frac{\dd}{\dd t}\right|_{t=0}\exp_G(tX)
 =\gamma_X'(0)=X.
\]
Różniczka jest zatem izomorfizmem. Twierdzenie o funkcji
odwrotnej z Rozdziału~\ref{ch:rozmaitosci} daje ostatni wniosek.
\end{proof}

\begin{przyklad}[Eksponenta macierzowa]
\label{prz:matrix-exp-lie}
\index{eksponenta macierzowa}
Dla $X\in M_n(\R)$ lub $M_n(\C)$ szereg
\[
 E_X(t)=\sum_{j=0}^{\infty}\frac{t^jX^j}{j!}
\]
jest zbieżny bezwzględnie i jednostajnie z pochodnymi
na każdym zwartym przedziale $t$: wybieramy normę
operatorową macierzy, w której norma składnika nie
przekracza $|t|^j\|X\|^j/j!$. Różniczkowanie wyraz po
wyrazie daje $E_X'(t)=E_X(t)X$, $E_X(0)=I$.
Macierz $X$ komutuje z każdym $E_X(t)$, więc
pochodna $E_X(t)E_X(-t)$ jest zerowa; jej wartość
w zerze to $I$. Zatem $E_X(t)^{-1}=E_X(-t)$
i cała krzywa leży w $\mathrm{GL}(n,\mathbb K)$.
Spełnia równanie pola $X^L_A=AX$ na
$\mathrm{GL}(n,\mathbb K)$; z jednoznaczności wynika
\[
 \exp_{\mathrm{GL}(n,\mathbb K)}(X)
 =E_X(1)=\sum_{j=0}^{\infty}\frac{X^j}{j!}.
\]
Jeżeli $X^T=-X$, to
$E_X(t)^T=E_{-X}(t)=E_X(t)^{-1}$;
ciągłość wyznacznika od $t=0$ pokazuje
$E_X(t)\in\mathrm{SO}(n)$.
Jeżeli
$X=\begin{psmallmatrix}i\alpha&z\\-\bar z&-i\alpha\end{psmallmatrix}
\in\mathfrak{su}(2)$, to
$X^2=-(\alpha^2+|z|^2)I$. Dla
$r=\sqrt{\alpha^2+|z|^2}>0$ szereg daje
$E_X(t)=\cos(rt)I+\frac{\sin(rt)}{r}X$;
przy $r=0$ jest to $I$. Ponieważ $X^*=-X$,
macierz $E_X(t)$ jest unitarna, a
$\det E_X(t)=\cos^2(rt)+\sin^2(rt)=1$.
To wyprowadza jej przynależność do $\mathrm{SU}(2)$
bez odwołania do ogólnego wzoru na wyznacznik
eksponenty.

Dla $J=\begin{psmallmatrix}0&-1\\1&0\end{psmallmatrix}$
zachodzi $J^2=-I$, więc rozdzielenie parzystych i
nieparzystych potęg daje
\[
 \exp(\theta J)=
 \begin{pmatrix}\cos\theta&-\sin\theta\\
 \sin\theta&\cos\theta\end{pmatrix}.
\]
Torus w modelu $(S^1)^k$ ma z kolei eksponentę
$(\theta_1,\ldots,\theta_k)\mapsto
(e^{i\theta_1},\ldots,e^{i\theta_k})$ i jądro
$(2\pi\Z)^k$; współrzędne $\theta$ są tutaj
$2\pi$ razy współrzędnymi $v$ z definicji ilorazowej.
Lokalna odwracalność eksponenty nie oznacza więc
globalnej różnowartościowości. W $\mathrm{GL}(n,\R)$
eksponenta nie pokrywa nawet całej grupy: każdy
$\exp X$ leży na drodze $t\mapsto\exp(tX)$ od $I$,
więc ma dodatni wyznacznik.
\end{przyklad}

\section{Sprzężenie i działanie adjungowane}
\label{sec:lie-adjoint}

\begin{definicja}[Działanie adjungowane]
\label{def:lie-ad}
\index{działanie adjungowane}
Dla $g\in G$ \emph{sprzężenie} przez $g$ to automorfizm
$c_g:G\to G$, $c_g(h)=ghg^{-1}$. Jego różniczkę przy $e$
oznaczamy
\[
 \operatorname{Ad}_g:=(\dd c_g)_e:
 \mathfrak g\longrightarrow\mathfrak g.
\]
Ponieważ $c_g\circ c_h=c_{gh}$, otrzymujemy
$\operatorname{Ad}_{gh}=\operatorname{Ad}_g\operatorname{Ad}_h$.
Jest to liniowy sposób zapisania, jak zmiana punktu odniesienia
przez sprzężenie zmienia prędkości przy jedności.
W grupie macierzy $\operatorname{Ad}_gX=gXg^{-1}$,
co wynika przez zróżniczkowanie
$g\exp_G(tX)g^{-1}$ w zerze.
\end{definicja}

\begin{twierdzenie}[Różniczka działania adjungowanego]
\label{thm:lie-ad-bracket}
Odwzorowanie $g\mapsto\operatorname{Ad}_g$ jest gładkim
homomorfizmem $G\to\mathrm{GL}(\mathfrak g)$.
Dla $X,Y\in\mathfrak g$ zachodzi
\begin{equation}\label{eq:lie-ad-bracket}
 \left.\frac{\dd}{\dd t}\right|_{t=0}
 \operatorname{Ad}_{\exp_G(tX)}Y=[X,Y]_{\mathfrak g}.
\end{equation}
Ponadto $\operatorname{Ad}_g[X,Y]_{\mathfrak g}
=[\operatorname{Ad}_gX,\operatorname{Ad}_gY]_{\mathfrak g}$.
\end{twierdzenie}
\begin{proof}
Gładkość sprzężenia jako odwzorowania $(g,h)\mapsto ghg^{-1}$
oraz gładkość jego różniczki dają gładkość $g\mapsto
\operatorname{Ad}_g$. Reguła łańcucha i
$c_g\circ c_h=c_{gh}$ dają własność homomorfizmu
i odwrotność $\operatorname{Ad}_{g^{-1}}$.
Ponieważ $c_g\circ L_h=L_{c_g(h)}\circ c_g$,
przeniesienie pola lewostronnego przez $c_g$ jest
ponownie lewostronne, z wartością przy $e$ równą
$\operatorname{Ad}_gY$:
\begin{equation}\label{eq:ad-pole}
 (c_g)_*Y^L=(\operatorname{Ad}_gY)^L.
\end{equation}

Ustalmy $X$ i połóżmy $a_t=\exp_G(tX)$.
Rodzina $c_{a_t}$ jest przepływem pola
$V=X^R-X^L$, gdzie $X^R_h=(\dd R_h)_eX$.
Istotnie, różniczkowanie $a_t h a_{-t}$ przy $t=0$
daje $(\dd R_h)_eX-(\dd L_h)_eX$, a prawo grupowe
tej rodziny pozwala zastosować jednoznaczność przepływu.
Przepływem $X^R$ jest $h\mapsto a_t h$, a przepływem
$Y^L$ jest $h\mapsto h\exp_G(tY)$. Translacje lewe
i prawe komutują, zatem z
Stwierdzenia~\ref{prop:inwarianty} wynika
$[X^R,Y^L]=0$. Korzystając z konwencji pochodnej
Liego z Twierdzenia~\ref{thm:lie-nawias}, mamy
$(c_{a_t})_*Y^L=(c_{a_{-t}})^*Y^L$,
więc pochodna pola przeniesionego przez przepływ $V$
wynosi $-[V,Y^L]$. Dlatego
\[
 \left.\frac{\dd}{\dd t}\right|_{t=0}
 (c_{a_t})_*Y^L
 =-[X^R-X^L,Y^L]=[X^L,Y^L].
\]
Ocena przy $e$ i~\eqref{eq:ad-pole} dają~\eqref{eq:lie-ad-bracket}.
Na koniec naturalność nawiasu względem $c_g$,
użyta już w dowodzie Twierdzenia~\ref{thm:lie-algebra-fields},
oraz~\eqref{eq:ad-pole} dają ostatnią równość.
\end{proof}

\begin{przyklad}[Macierze i torus]
\label{prz:ad-matrix-torus}
Dla grupy macierzy możemy niezależnie sprawdzić znak
w~\eqref{eq:lie-ad-bracket}:
\[
 \left.\frac{\dd}{\dd t}\right|_0
 e^{tX}Ye^{-tX}=XY-YX.
\]
W grupie przemiennej każde sprzężenie jest identycznością.
Stąd $\operatorname{Ad}$ torusa jest trywialne, a jej
różniczka i nawias algebry Liego znikają.
W teorii Cherna--Weila z Rozdziału~\ref{ch:chern-weil}
warunek niezmienniczości wielomianu oznacza właśnie
niezmienniczość względem takich przekształceń
$\operatorname{Ad}_g$.
\end{przyklad}

\section{Od SU(2) do obrotów przestrzeni}
\label{sec:su2-so3-lie}

W Przykładzie~\ref{prz:spin3} grupa
$\operatorname{Spin}(3)$ została utożsamiona z jednostkowymi
kwaternionami. Tutaj połączymy ten opis z macierzami
$\mathrm{SU}(2)$ i obliczymy samo nakrycie.

Zapiszmy kwaternion jako $q=a+b\mathbf j$, gdzie $a,b\in\C$,
$\mathbf j^2=-1$ i $\mathbf jz=\bar z\mathbf j$.
Iloczyn ma wtedy postać
\[
 (a+b\mathbf j)(c+d\mathbf j)
 =(ac-b\bar d)+(ad+b\bar c)\mathbf j.
\]
Odwzorowanie $q\mapsto U(a,b)$ z~\eqref{eq:su2-macierz}
zachowuje ten iloczyn, co widać po pomnożeniu macierzy:
jej pierwszy wiersz ma współczynniki $ac-b\bar d$
i $ad+b\bar c$. Ponadto
$|q|^2=|a|^2+|b|^2$. Jest to więc gładki izomorfizm
grup jednostkowych kwaternionów z $\mathrm{SU}(2)$.

\begin{twierdzenie}[Jawne dwukrotne nakrycie]
\label{thm:su2-so3-cover-lie}
\index{nakrycie!$\mathrm{SU}(2)$ nad $\mathrm{SO}(3)$}
Niech $\operatorname{Im}\mathbb H$ będzie przestrzenią
kwaternionów o zerowej części rzeczywistej, utożsamioną
z $\R^3$ przez bazę $(\mathbf i,\mathbf j,\mathbf k)$.
Wtedy
\[
 \rho:\mathrm{SU}(2)\longrightarrow\mathrm{SO}(3),
 \qquad \rho(q)(v)=qvq^{-1},\quad
 v\in\operatorname{Im}\mathbb H,
\]
jest gładkim surjektywnym homomorfizmem o jądrze
$\{1,-1\}$. Każdy punkt $\mathrm{SO}(3)$ ma otoczenie
równomiernie pokryte przez dwa otwarte zbiory w
$\mathrm{SU}(2)$; zatem $\rho$ jest dwukrotnym nakryciem.
\end{twierdzenie}
\begin{proof}
Kwaternionowy iloczyn spełnia $|qr|=|q||r|$ oraz
$\Re(ab)=\Re(ba)$, co wynika z reguł mnożenia bazy.
Dla $|q|=1$ sprzężenie zachowuje więc część rzeczywistą
i normę, a zatem ogranicza się do izometrii liniowej
$\operatorname{Im}\mathbb H$. Zależność od $q$ jest
gładka, bo $q^{-1}=\bar q$ na sferze jednostkowej.
Grupa $S^3$ jest spójna drogowo: każdy jej element
można zapisać jako $\cos t+u\sin t$ dla jednostkowego
urojonego $u$ (przy $q=\pm1$ wybór $u$ jest dowolny).
Wyznacznik izometrii $\rho(q)$ jest więc stale równy
$+1$, jak dla $q=1$. Łączność mnożenia daje
$\rho(qr)=\rho(q)\rho(r)$.

Dla urojonych kwaternionów $u,v$ mamy
$uv=-\langle u,v\rangle+u\times v$.
Niech $|u|=1$ i
$q=\cos(\theta/2)+u\sin(\theta/2)$.
Ponieważ $uv-vu=2u\times v$ i
$uvu=v-2\langle u,v\rangle u$, rozwinięcie
$(c+su)v(c-su)$ dla
$c=\cos(\theta/2)$, $s=\sin(\theta/2)$ daje
\begin{align*}
 \rho(q)v
 &=c^2v+cs(uv-vu)-s^2uvu\\
 &=\cos\theta\,v+
   \sin\theta\,(u\times v)+
   (1-\cos\theta)\langle u,v\rangle u.
\end{align*}
Jest to obrót o kąt $\theta$ wokół osi $u$.
Każda macierz $A\in\mathrm{SO}(3)$ ma oś stałą:
\[
 \det(I-A)=\det(I-A^T)
 =\det(A^T)\det(A-I)
 =\det(A-I)=-\det(I-A),
\]
więc $\det(I-A)=0$. Wybierzmy jednostkowy
$u\in\ker(I-A)$. Na $u^\perp$ macierz $A$ jest
izometrią o wyznaczniku $+1$. Orientujemy tę płaszczyznę
tak, by $(u,\text{dodatnia baza }u^\perp)$ była dodatnia
w $\R^3$; wtedy $A|_{u^\perp}$ jest obrotem pewnego
kąta $\theta$. Powyższy wzór daje kwaternion,
który realizuje dokładnie $A$; $\rho$ jest surjektywna.

Jeśli $\rho(q)$ jest identycznością, to $q$ komutuje
ze wszystkimi urojonymi $v$. Po zapisaniu $q=a+u$
z $a\in\R$, $u\in\operatorname{Im}\mathbb H$,
otrzymujemy $0=qv-vq=2u\times v$ dla każdego $v$.
Stąd $u=0$ i, przez $|q|=1$, $q=\pm1$.

Pozostaje własność nakrycia, a nie tylko policzenie
dwóch przeciwobrazów. Dla $u\in\operatorname{Im}\mathbb H$
krzywa $q(t)=\cos(t|u|)+(u/|u|)\sin(t|u|)$
(dla $u=0$ stała) ma pochodną $u$ w zerze.
Różniczkując sprzężenie, dostajemy
\[
 (\dd\rho)_1(u)(v)=uv-vu=2u\times v.
\]
Jest to izomorfizm $\operatorname{Im}\mathbb H\to
\mathfrak{so}(3)$: jest różnowartościowy, bo
$u\times v=0$ dla wszystkich $v$ wymusza $u=0$,
a obie przestrzenie mają wymiar trzy.
Ponieważ $\rho$ jest homomorfizmem, translacje
przenoszą tę różniczkę na dowolny punkt. Twierdzenie
o funkcji odwrotnej daje więc otoczenie $W$ jedności,
na którym $\rho$ jest dyfeomorfizmem. Zmniejszmy $W$
tak, aby $W\cap(-W)=\varnothing$, i połóżmy
$V=\rho(W)$. Każdy punkt $V$ ma dokładnie
dwa przeciwobrazy, zatem
$\rho^{-1}(V)=W\sqcup(-W)$ i oba ograniczenia są
dyfeomorfizmami na $V$. Translacja przenosi tę
konstrukcję na otoczenie dowolnego obrotu. To jest
definicja dwukrotnego nakrycia z~\ref{def:nakrycie}.
\end{proof}

W szczególności $\mathrm{SU}(2)\cong\operatorname{Spin}(3)$
jako grupy Liego, a powyższe $\rho$ jest tym samym
nakryciem co w Twierdzeniu~\ref{thm:spin-nakrycie}.
Rzeczywiście w algebrze Clifforda z
Przykładu~\ref{prz:spin3} element
$\omega=e_1e_2e_3$ jest centralny, a
$\mathcal J(v)=-v\omega$ wysyła bazę $(e_1,e_2,e_3)$ na
$(I,J,K)$ przestrzeni urojonych kwaternionów.
Ponieważ $\omega$ komutuje z $s\in\operatorname{Spin}(3)$,
otrzymujemy $\mathcal J(svs^{-1})=s\mathcal J(v)s^{-1}$, czyli obie
konstrukcje nakrycia są zgodne pod tym utożsamieniem.
Jego różniczka jest izomorfizmem algebr Liego
$\mathfrak{su}(2)\cong\mathfrak{so}(3)$, lecz same grupy
Liego nie są izomorficzne. Istotnie,
$\mathrm{SU}(2)\cong S^3$ jest jednospójna na mocy
Przykładu~\ref{prz:pi1-sfera}. Droga
$q(t)=\cos(\pi t)+\mathbf i\sin(\pi t)$, $0\leq t\leq1$,
łączy $1$ z $-1$, a jej obraz przez $\rho$ jest pętlą
w $\mathrm{SO}(3)$. Gdyby ta pętla była ściągalna,
Twierdzenie~\ref{thm:podnoszenie-nakrycie} pozwoliłoby
podnieść jej homotopię do homotopii dróg o stałych
końcach; byłoby to sprzeczne z faktem, że jej podniesienie
zaczyna się w $1$, a kończy w $-1$.
Zatem $\pi_1(\mathrm{SO}(3))$ nie jest trywialna.

\par\smallskip
{\footnotesize\noindent\emph{Dalsza lektura:}
\href{https://archive.org/details/theotransformation01liesrich}
{S.\ Lie i F.\ Engel, \emph{Theorie der Transformationsgruppen},
t.~I (1888)}, oryginalne ujęcie grup przekształceń i pól
infinitesymalnych;
\href{https://books.google.com/books/about/Theory_of_Lie_Groups_I.html?id=bm4GAQAAIAAJ}
{C.\ Chevalley, \emph{Theory of Lie Groups}, t.~I (1946)},
systematyczne ujęcie grup gładkich i ich algebr.\par}
